%=======================================================================
% resumen.tex
%=======================================================================
Chile alberga 528 especies de aves nativas, 12 de ellas endémicas, y 71 de
las 147 clasificadas por el Reglamento de Clasificación de Especies están
amenazadas \citep{mma2026rce}. Su monitoreo depende de especialistas y de
trabajo de campo costoso, mientras la ciencia ciudadana y los satélites
acumulan datos que apenas se aprovechan. Este trabajo desarrolla y evalúa un
sistema de aprendizaje profundo que identifica especies de aves chilenas a
partir de audio y/o imagen mediante fusión multimodal, y que modela su
distribución espaciotemporal a partir de avistamientos, clima e índices de
vegetación.

Se integraron en un repositorio reproducible grabaciones de Xeno-canto
(\num{2508}), fotografías de iNaturalist vía GBIF (\num{7060}),
\num{100270} avistamientos de eBird e iNaturalist y covariables ambientales
de Google Earth Engine, para 55 especies de Chile continental con foco en
las amenazadas y endémicas. En audio se ajustan una CNN (EfficientNet-B0) y
el \textit{Audio Spectrogram Transformer} sobre espectrogramas log-mel; en
imagen, EfficientNet-B3 con aumentación en tiempo de prueba. Ambas ramas se
combinan por \textit{late fusion} y se contrastan con clasificadores
unimodales, con \textit{frozen embeddings} y con un modelo de IA generativa
multimodal. En distribución se comparan MaxEnt, Random Forest, gradient
boosting y una red neuronal multiespecie, con partición aleatoria y
validación espacial por bloques.

La fusión identifica 53 especies con 0,945 de \textit{accuracy} (intervalo
de confianza al \SI{95}{\percent} de 0,924 a 0,959), por sobre el audio solo
(0,865), la imagen sola (0,816), el modelo generativo (0,817) y la meta del
\SI{85}{\percent}. Random Forest supera a MaxEnt en AUC-ROC con partición
aleatoria (0,884 frente a 0,863, $p < 0{,}001$), pero bajo validación
espacial todos los modelos convergen a 0,80--0,82; la red neuronal no supera
a MaxEnt y los \textit{embeddings} satelitales no aportan.

El sistema cumple sus objetivos y deja tres aportes: la evidencia de que la
fusión audio-imagen con datos no emparejados de ciencia ciudadana supera a
cada modalidad; un resultado negativo bien fundamentado sobre el aprendizaje
profundo en modelos de distribución tabulares; y un repositorio reproducible
con un prototipo integrado, extensible a otras especies y regiones.
